\documentclass[11pt]{article}
\usepackage{fontspec,amsmath,amssymb,amsthm,hyperref}
\usepackage[margin=1in]{geometry}
\setlength{\emergencystretch}{3em}
\title{Uniform Euler mesh convergence in the actual L\ensuremath{{}^{2}} space}
\author{LangevinLean mathematical exposition}
\date{}
\begin{document}\maketitle
\noindent Source module: \texttt{AverageLMC/LaplacianEulerLimit.lean}. Compile this proof document with XeLaTeX or LuaLaTeX. The accompanying HTML pairs the mathematical argument with the exact Lean statements and records the scope of the formalization.


The existing additive integral equation is used directly in a Banach space. Matching-noise Euler convergence is strengthened from its terminal value to every mesh point, giving the uniform drift moments required to sum weak errors.



Let \(E_k^{(n)}\) be the matching-noise Euler chain with mesh \(h_n=T/(n+1)\). The integral equation implies simultaneous mesh convergence:


\[\max_{0\le k\le n+1}\left\|E_k^{(n)}-Y(kh_n)\right\|
\longrightarrow0.\]


For every uniformly continuous observation \(F\) into a Banach space,


\[\sum_{k=0}^{n}h_nF\bigl(E_k^{(n)}\bigr)
\longrightarrow\int_0^T F(Y(t))\,dt.\]
\section{Estimate every grid point with one drift modulus}


Restrict the existing Euler drift-modulus estimate to the shorter interval ending at kh. Monotonicity of kh and the exponential factor gives the same upper bound T ε exp(LT) at every point k\ensuremath{\le}n.



\noindent Lean declarations: \texttt{additiveEulerChain\_grid\_error\_le\_of\_drift\_modulus}.

\section{Let the compact-time modulus shrink}


The drift along the continuous integral solution is uniformly continuous on [0,T]. Choose a drift tolerance from the requested path error, then refine the mesh below its continuity radius. The preceding estimate proves simultaneous convergence of all Euler grid points.



\noindent Lean declarations: \texttt{additiveEulerChain\_grid\_uniform\_tendsto}.

\section{Derive uniform state and score moments}


The true continuous path has bounded norm on the compact time interval. Eventually all Euler grid points are within distance one of this path. Lipschitz growth then bounds every drift value by L times this common state radius plus ‖g(0)‖. In the L\ensuremath{{}^{2}} application these are actual uniform second-moment bounds.



\noindent Lean declarations: \texttt{additiveEulerChain\_grid\_eventually\_bounded}, \texttt{additiveEulerChain\_drift\_grid\_eventually\_bounded}.

\section{Prove the uniform left-sum estimate}


On each partition cell, integrate the difference between the left-endpoint value and the actual continuous integrand. A modulus ε bounds its norm by hε. The adjacent-interval identity and the norm of a finite sum give total error at most Tε. Compact-time uniform continuity then proves convergence of the uniform Riemann sums.



\noindent Lean declarations: \texttt{uniform\_left\_sum\_error\_le\_of\_modulus}, \texttt{uniform\_left\_sum\_tendsto}.

\section{Pass observations of Euler states to time integrals}


Uniform continuity of an observation sends the simultaneous Euler grid error to one common observation error. Sum this error over the partition, and combine it with the continuous path's left Riemann-sum limit. No differentiability of the Brownian path is required.



\noindent Lean declarations: \texttt{additiveEulerChain\_test\_sum\_tendsto}.

\end{document}
